% Mathématiques 4e — cours V3
% Grandeurs, espace et volumes

\section{Objectifs}

\begin{itemize}
\item Se repérer dans un pavé droit à l'aide de trois coordonnées.
\item Lire et produire une perspective simple d'un solide.
\item Reconnaître et exploiter des patrons de pyramides et de cônes.
\item Calculer le volume d'une pyramide.
\item Calculer le volume d'un cône de révolution.
\item Utiliser des grandeurs produits et quotients : vitesse, débit, concentration, masse volumique.
\item Convertir des unités composées.
\item Contrôler la cohérence des unités dans une formule.
\item Résoudre des problèmes combinant géométrie, volumes et grandeurs composées.
\end{itemize}

\section{1. Repérage dans l'espace}

Dans un pavé droit, on peut repérer un point à l'aide de trois coordonnées.

On utilise par exemple :
\[
(x;y;z).
\]

Selon le contexte, ces coordonnées peuvent être appelées :
\begin{itemize}
\item abscisse ;
\item ordonnée ;
\item altitude.
\end{itemize}

\coursefigure{/images/mathematiques/4eme/espace-repere-3d.svg}{Pavé droit muni de trois directions de repérage.}{Le repérage dans l'espace ajoute une troisième coordonnée au repérage dans le plan.}

\section{2. Lire trois coordonnées}

Si :
\[
P(3;2;5),
\]
cela signifie que le point est repéré par trois déplacements selon les directions choisies.

Il faut toujours vérifier :
\begin{itemize}
\item l'origine ;
\item l'ordre des axes ;
\item les unités ;
\item les graduations.
\end{itemize}

\section{3. Perspective}

Une perspective représente un solide de l'espace sur un support plan.

Elle ne conserve pas nécessairement toutes les longueurs et tous les angles.

Les arêtes cachées peuvent être dessinées en pointillés.

\begin{attention}
Mesurer directement sur une perspective ne donne pas nécessairement une longueur réelle du solide.
\end{attention}

\section{4. Pyramide}

Une pyramide possède :
\begin{itemize}
\item une base polygonale ;
\item un sommet qui n'appartient pas au plan de la base ;
\item des faces latérales triangulaires.
\end{itemize}

La hauteur de la pyramide est la distance perpendiculaire du sommet au plan de la base.

\section{5. Cône de révolution}

Un cône de révolution possède :
\begin{itemize}
\item une base circulaire ;
\item un sommet ;
\item une hauteur perpendiculaire au plan de la base.
\end{itemize}

Le rayon de la base intervient dans son aire :
\[
A_{\text{base}}=\pi r^2.
\]

\coursefigure{/images/mathematiques/4eme/espace-pyramide-cone.svg}{Pyramide et cône avec base et hauteur représentées.}{Le volume d'une pyramide ou d'un cône vaut un tiers du produit de l'aire de base par la hauteur.}

\section{6. Volume d'une pyramide}

\begin{propriete}
Pour une pyramide d'aire de base $B$ et de hauteur $h$ :
\[
\boxed{V=\dfrac{B\times h}{3}}.
\]
\end{propriete}

\begin{exemple}
Une pyramide possède une base rectangulaire de :
\[
6\ \text{cm}\times4\ \text{cm}
\]
et une hauteur de :
\[
9\ \text{cm}.
\]

L'aire de base vaut :
\[
B=6\times4=24\ \text{cm}^2.
\]

Donc :
\[
V=\dfrac{24\times9}{3}=72\ \text{cm}^3.
\]
\end{exemple}

\section{7. Volume d'un cône}

Pour un cône de rayon $r$ et hauteur $h$ :
\[
B=\pi r^2.
\]

Donc :
\[
\boxed{
V=\dfrac{\pi r^2h}{3}
}.
\]

\section{8. Exemple développé : cône}

Un cône possède :
\[
r=3\ \text{cm}
\]
et :
\[
h=8\ \text{cm}.
\]

Alors :
\[
V=\dfrac{\pi\times3^2\times8}{3}.
\]

Donc :
\[
V=24\pi\ \text{cm}^3.
\]

Valeur approchée :
\[
V\approx75{,}4\ \text{cm}^3.
\]

\section{9. Comparer cylindre et cône}

Un cylindre et un cône possèdent la même base et la même hauteur.

Le cylindre a pour volume :
\[
V_{\text{cyl}}=Bh.
\]

Le cône :
\[
V_{\text{cône}}=\dfrac{Bh}{3}.
\]

Donc :
\[
\boxed{
V_{\text{cône}}=\dfrac13V_{\text{cyl}}
}.
\]

\section{10. Patron d'une pyramide}

Un patron de pyramide comporte :
\begin{itemize}
\item la base ;
\item une face triangulaire attachée à chaque côté de la base.
\end{itemize}

Les longueurs qui se rejoignent au pliage doivent être compatibles.

Un dessin qui contient les bonnes formes mais des longueurs incompatibles n'est pas un patron valable.

\section{11. Patron d'un cône}

Le patron d'un cône comporte :
\begin{itemize}
\item un disque pour la base ;
\item un secteur de disque pour la surface latérale.
\end{itemize}

L'arc du secteur doit avoir la même longueur que le périmètre de la base.

Cette propriété permet de contrôler un patron.

\section{12. Grandeur quotient}

Certaines grandeurs sont obtenues en divisant deux grandeurs.

Exemple :
\[
v=\dfrac dt.
\]

La vitesse est une distance divisée par une durée.

Si :
\[
d=150\ \text{km}
\]
et :
\[
t=2\ \text{h},
\]
alors :
\[
v=75\ \text{km/h}.
\]

\section{13. Vitesse}

Les trois relations sont :
\[
\boxed{v=\dfrac dt},
\]
\[
\boxed{d=vt},
\]
\[
\boxed{t=\dfrac dv}.
\]

Il faut choisir celle qui correspond à la grandeur recherchée.

\section{14. Conversion de vitesse}

Pour convertir :
\[
72\ \text{km/h}
\]
en :
\[
\text{m/s},
\]
on transforme :
\[
72\ \text{km}=72\,000\ \text{m},
\]
et :
\[
1\ \text{h}=3600\ \text{s}.
\]

Donc :
\[
72\ \text{km/h}
=
\dfrac{72\,000}{3600}\ \text{m/s}
=
20\ \text{m/s}.
\]

\section{15. Débit}

Un débit peut être défini par :
\[
\boxed{q=\dfrac{V}{t}}.
\]

Si :
\[
120\ \text{L}
\]
s'écoulent en :
\[
8\ \text{min},
\]
alors :
\[
q=15\ \text{L/min}.
\]

\section{16. Concentration simple}

Dans certaines situations, une concentration est un quotient :
\[
c=\dfrac{\text{quantité dissoute}}{\text{volume de solution}}.
\]

Par exemple :
\[
12\ \text{g}
\]
de soluté dans :
\[
0{,}5\ \text{L}
\]
donnent :
\[
c=\dfrac{12}{0{,}5}=24\ \text{g/L}.
\]

Le contexte scientifique donne le sens précis de la grandeur.

\section{17. Masse volumique}

La masse volumique est :
\[
\boxed{\rho=\dfrac mV}.
\]

Si un objet a une masse :
\[
540\ \text{g}
\]
et un volume :
\[
200\ \text{cm}^3,
\]
alors :
\[
\rho=2{,}7\ \text{g/cm}^3.
\]

\section{18. Grandeur produit}

Certaines grandeurs sont obtenues par multiplication.

Par exemple :
\[
d=vt.
\]

La distance est ici le produit d'une vitesse par une durée compatible.

Le contrôle d'unités donne :
\[
\text{km/h}\times\text{h}=\text{km}.
\]

\section{19. Cohérence des unités}

Considérons :
\[
v=60\ \text{km/h},
\]
et :
\[
t=30\ \text{min}.
\]

On ne peut pas écrire directement :
\[
d=60\times30.
\]

Il faut convertir :
\[
30\ \text{min}=0{,}5\ \text{h}.
\]

Donc :
\[
d=60\times0{,}5=30\ \text{km}.
\]

\section{20. Unités composées}

Des écritures comme :
\[
\text{km/h},
\quad
\text{L/min},
\quad
\text{g/L},
\quad
\text{kg/m}^3
\]
sont des unités composées.

Le trait de fraction ou le symbole « par » doit être interprété comme un quotient.

\section{21. Convertir une masse volumique}

On sait :
\[
1\ \text{g/cm}^3=1000\ \text{kg/m}^3.
\]

Ainsi :
\[
2{,}7\ \text{g/cm}^3
=
2700\ \text{kg/m}^3.
\]

Les conversions de grandeurs composées demandent de convertir le numérateur et le dénominateur de manière cohérente.

\section{22. Problème mêlant volume et capacité}

Un réservoir a la forme d'un pavé de dimensions :
\[
80\ \text{cm},
\quad
50\ \text{cm},
\quad
40\ \text{cm}.
\]

Son volume vaut :
\[
80\times50\times40=160\,000\ \text{cm}^3.
\]

Or :
\[
1000\ \text{cm}^3=1\ \text{L}.
\]

Donc :
\[
V=160\ \text{L}.
\]

\section{23. Problème de cône et débit}

Un récipient conique de rayon :
\[
6\ \text{cm}
\]
et hauteur :
\[
15\ \text{cm}
\]
est rempli à débit constant.

Son volume est :
\[
V=\dfrac{\pi\times6^2\times15}{3}
=
180\pi\ \text{cm}^3.
\]

Comme :
\[
1\ \text{cm}^3=1\ \text{mL},
\]
le volume est environ :
\[
565\ \text{mL}.
\]

À un débit :
\[
50\ \text{mL/s},
\]
le temps vaut environ :
\[
565\div50\approx11{,}3\ \text{s}.
\]

\section{24. Méthode volume}

\begin{methode}
\begin{enumerate}
\item identifier le solide ;
\item repérer la base et la hauteur perpendiculaire ;
\item calculer l'aire de base ;
\item choisir la formule de volume ;
\item harmoniser les unités ;
\item calculer ;
\item donner l'unité cubique ou convertir en capacité.
\end{enumerate}
\end{methode}

\section{25. Méthode grandeur composée}

\begin{methode}
\begin{enumerate}
\item identifier les grandeurs connues ;
\item écrire la relation entre elles ;
\item vérifier la compatibilité des unités ;
\item convertir si nécessaire ;
\item effectuer le calcul ;
\item contrôler l'unité finale et l'ordre de grandeur.
\end{enumerate}
\end{methode}

\section{26. Erreurs fréquentes}

\begin{itemize}
\item Oublier le facteur $\dfrac13$ pour une pyramide ou un cône.
\item Utiliser la longueur d'une arête oblique comme hauteur.
\item Confondre aire de base et périmètre de base.
\item Mélanger minutes et heures dans une formule de vitesse.
\item Confondre $\text{cm}^2$ et $\text{cm}^3$.
\item Convertir une unité composée en ne modifiant qu'une partie.
\item Lire les trois coordonnées dans le mauvais ordre.
\end{itemize}

\section{À retenir}

\begin{aretenir}
Pour une pyramide ou un cône :
\[
\boxed{
V=\dfrac{B\times h}{3}
}.
\]

Pour un cône :
\[
\boxed{
V=\dfrac{\pi r^2h}{3}
}.
\]

Les grandeurs composées se lisent à travers leurs unités :
\[
\boxed{
v=\dfrac dt,\qquad
q=\dfrac Vt,\qquad
\rho=\dfrac mV
}.
\]

Avant tout calcul, les unités doivent être rendues compatibles.
\end{aretenir}
